\documentclass[11pt,a4paper]{article}

% --- Encoding & language ---------------------------------------------------
\usepackage[utf8]{inputenc}
\usepackage[T1]{fontenc}
\usepackage[french]{babel}
\usepackage{lmodern}
\usepackage{microtype}

% --- Layout ----------------------------------------------------------------
\usepackage[a4paper, top=2.4cm, bottom=2.4cm, left=2.5cm, right=2.5cm]{geometry}
\usepackage{setspace}
\onehalfspacing

% --- Visuals ---------------------------------------------------------------
\usepackage{xcolor}
\definecolor{kvink}{RGB}{12,18,28}
\definecolor{kvgold}{RGB}{196,142,42}
\definecolor{kvalert}{RGB}{226,58,58}
\definecolor{kvgreen}{RGB}{63,163,108}
\definecolor{kvgrey}{RGB}{120,128,140}
\definecolor{kvbg}{RGB}{246,246,242}

\usepackage{titlesec}
\titleformat{\section}
  {\color{kvink}\Large\bfseries\sffamily}
  {\thesection.}{0.6em}{}
\titleformat{\subsection}
  {\color{kvink}\large\bfseries\sffamily}
  {\thesubsection}{0.6em}{}
\titleformat{\subsubsection}
  {\color{kvgold}\normalsize\bfseries\sffamily}
  {}{0pt}{}

% --- Tables & graphics -----------------------------------------------------
\usepackage{booktabs}
\usepackage{array}
\usepackage{tabularx}
\usepackage{multirow}
\renewcommand{\arraystretch}{1.25}

\usepackage{enumitem}
\setlist{leftmargin=1.4em, itemsep=2pt, topsep=3pt}

% --- Header & footer -------------------------------------------------------
\usepackage{fancyhdr}
\pagestyle{fancy}
\fancyhf{}
\fancyhead[L]{\sffamily\color{kvink}\textbf{KOVER.IA}}
\fancyhead[R]{\sffamily\color{kvgrey}\small Hackathon AI Paris~2026}
\fancyfoot[C]{\sffamily\color{kvgrey}\small \thepage}
\renewcommand{\headrulewidth}{0.4pt}
\renewcommand{\footrulewidth}{0pt}

% --- Boxes -----------------------------------------------------------------
\usepackage{tcolorbox}
\tcbuselibrary{skins,breakable}
\newtcolorbox{insight}[1][]{
  enhanced, breakable, colback=kvbg, colframe=kvink,
  boxrule=0.5pt, arc=2pt, left=10pt, right=10pt, top=8pt, bottom=8pt,
  fontupper=\sffamily\small, #1
}
\newtcolorbox{quote_kv}[1][]{
  enhanced, colback=white, colframe=kvgold,
  boxrule=1pt, arc=0pt, left=14pt, right=14pt, top=10pt, bottom=10pt,
  borderline west={3pt}{0pt}{kvgold},
  fontupper=\itshape, #1
}

% --- Hyperref --------------------------------------------------------------
\usepackage[hidelinks]{hyperref}

% ===========================================================================
\begin{document}

% --- Title page -----------------------------------------------------------
\begin{titlepage}
\centering
\vspace*{2.5cm}

{\sffamily\Huge\bfseries\color{kvink} KOVER.IA}\\[0.2cm]
{\sffamily\Large\color{kvgold} \rule{4cm}{1pt}}\\[0.4cm]

{\sffamily\LARGE\bfseries Détection en temps réel \\[0.2cm] du blanchiment crypto par IA}\\[1.0cm]

{\sffamily\large\color{kvgrey} Pipeline multi-modèle : GNN + RNN + LLM\\[0.1cm]
+ moteur d'interception de flashloan attacks}\\[3.2cm]

\begin{tcolorbox}[
  enhanced, colback=kvbg, colframe=kvink,
  boxrule=0.5pt, arc=2pt, width=12cm,
  left=12pt, right=12pt, top=10pt, bottom=10pt,
  fontupper=\sffamily
]
\centering
\textbf{Rapport projet --- Hackathon AI Paris 2026}\\[0.4cm]
\small Soumis dans le cadre de la compétition \\
solutions IA appliquées à la sécurité DeFi.
\end{tcolorbox}

\vfill

{\sffamily\color{kvgrey}\small Avril 2026 \quad\textbar\quad Version 1.0 \quad\textbar\quad \texttt{kover.ia}}
\end{titlepage}

% --- TOC -------------------------------------------------------------------
\renewcommand{\contentsname}{Table des matières}
{\sffamily\tableofcontents}
\thispagestyle{fancy}
\newpage

% ===========================================================================
\section{Introduction}

L'écosystème de la finance décentralisée (DeFi) traite aujourd'hui plus de
\textbf{120~milliards de dollars} de valeur verrouillée à travers des
protocoles ouverts, sans intermédiaire, opérant 24~heures sur~24 sur des
blockchains publiques. Cette ouverture, qui constitue précisément la
proposition de valeur du secteur, est aussi son talon d'Achille : selon
Chainalysis, plus de \textbf{2,2~milliards de dollars} ont été volés en
2024 par des attaques sur des protocoles DeFi, soit une augmentation de
21~\% par rapport à 2023. À ces vols directs s'ajoute le blanchiment qui
suit : les fonds dérobés transitent en quelques minutes par des mixeurs,
des bridges cross-chain et des plateformes d'échange centralisées (CEX),
souvent avant même que les équipes de sécurité ne soient alertées.

Les outils existants de surveillance anti-blanchiment (AML) sur la
blockchain sont massivement orientés vers l'analyse rétrospective :
ils étiquettent des adresses suspectes après les faits, alimentent des
listes de sanctions et produisent des rapports forensiques pour les
régulateurs. Ce qui manque cruellement à l'écosystème, c'est un système
\textbf{prédictif et temps réel} : un capteur qui voit un flux de
transactions en formation, le modélise comme un graphe vivant, et
émet une alerte sur la trajectoire probable des fonds avant qu'ils
n'atteignent un mixeur ou un dépôt CEX non régulé.

\textsc{KOVER.IA} a été conçu pour combler ce vide. Le projet combine
trois familles de modèles d'apprentissage automatique~--- un
\textbf{Graph Attention Network} entraîné sur le dataset Salam Ammari
(71~250 transactions Ethereum réelles), un \textbf{LSTM séquentiel}
qui prédit la prochaine destination d'un flux suspect, et un
\textbf{Large Language Model (Qwen 235B)} hébergé sur l'infrastructure
wafer-scale de Cerebras qui génère un rapport d'incident actionnable~---
en un pipeline orchestré qui consomme un flux de transactions live et
produit, en moins de cinq secondes, une décision motivée pour les
équipes de conformité ou les DAO de gouvernance.

À cette brique d'analyse de blanchiment s'ajoute une seconde capacité,
le module \textbf{Behavioral Flow Detection}, qui cible spécifiquement
les attaques par \emph{flashloan} : un sentinel surveille le mempool
public d'Ethereum, identifie les profils de transaction caractéristiques
d'un drainage imminent, et déclenche une riposte programmable
(\emph{halt}, alerte DAO, gel d'adresse) avant que la transaction
ne soit minée. Ces deux modules partagent la même infrastructure, la
même base de modèles, et la même interface de visualisation temps réel.

Ce rapport présente la problématique adressée, l'architecture technique
de la solution, le positionnement face aux concurrents établis, la
proposition de valeur unique de \textsc{KOVER.IA}, et le modèle
économique envisagé pour porter le projet vers une commercialisation
durable.

% ===========================================================================
\section{Problématique}

\subsection{Le coût caché du blanchiment crypto}

Lorsqu'une attaque réussit sur un protocole DeFi, l'attaquant n'a pas
gagné l'argent~: il l'a uniquement déplacé. La phase critique commence
alors~--- le blanchiment~--- et c'est là que la communauté perd presque
toujours la course. Les patterns observés sur les hacks récents
(Euler Finance, Ronin Bridge, Wormhole, Curve Finance) sont remarquablement
\textbf{stéréotypés}~:

\begin{enumerate}
  \item \textbf{Splitting initial} --- l'attaquant éclate les fonds sur
        4 à 12 wallets satellites en moins de deux minutes.
  \item \textbf{Mixing / peel chain} --- chaque wallet effectue 5 à 20
        transactions intermédiaires sur des mixeurs (Tornado Cash,
        Railgun) ou des séquences de wallets jetables.
  \item \textbf{Bridging cross-chain} --- les fonds transitent vers BSC,
        Polygon, Arbitrum via Wormhole, Stargate, Synapse.
  \item \textbf{Off-ramp} --- dépôt sur un CEX peu régulé pour une
        sortie en monnaie fiat.
\end{enumerate}

L'ensemble de cette séquence se déroule typiquement en
\textbf{15 à 90 minutes}. Les humains des équipes de sécurité, mêmes
expérimentés, sont structurellement incapables de réagir à cette
vélocité. Lorsqu'un analyste ouvre une investigation, les fonds sont
souvent déjà dans un mixeur ou sur un bridge, et la traçabilité on-chain
devient exponentiellement plus coûteuse.

\subsection{Limites des outils actuels}

Les solutions de surveillance AML existantes (Chainalysis Reactor,
TRM Labs, Elliptic Lens) reposent sur trois piliers historiques~:

\begin{itemize}
  \item \emph{Listes d'adresses étiquetées} maintenues manuellement par
        des équipes d'analystes (sanctions OFAC, mixeurs connus,
        wallets de hackers identifiés).
  \item \emph{Heuristiques de propagation} fondées sur des règles
        métier~: nombre de \emph{hops}, montants, fréquence.
  \item \emph{Tableaux de bord d'investigation} optimisés pour
        l'analyse rétrospective sur des fenêtres de jours ou semaines.
\end{itemize}

Ce paradigme produit des rapports utiles \emph{après} les incidents,
mais reste \textbf{aveugle au présent}~: il ne modélise pas le graphe
de transactions comme un système dynamique, n'apprend pas les
représentations latentes des wallets et de leurs voisinages, et ne
prédit jamais la prochaine étape probable d'un flux suspect.

\subsection{Le besoin d'un capteur prédictif temps réel}

La problématique adressée par \textsc{KOVER.IA} se formule ainsi~:

\begin{quote_kv}
Comment construire un système qui, en observant un flux de transactions
Ethereum live, détecte en moins de cinq secondes l'amorce d'une
opération de blanchiment, prédit sa destination probable, et émet un
rapport d'incident exploitable par une équipe de sécurité ou une DAO~?
\end{quote_kv}

Cette question soulève quatre sous-problèmes techniques distincts qu'il
a fallu résoudre simultanément~:

\begin{enumerate}
  \item \textbf{Représenter} chaque wallet par un vecteur appris,
        pas par une étiquette manuelle.
  \item \textbf{Capturer la dynamique séquentielle} d'un blanchiment
        (un acte n'a de sens que dans la séquence où il s'inscrit).
  \item \textbf{Produire un narratif actionnable} en langage naturel
        que les équipes peuvent injecter dans un workflow de réponse.
  \item \textbf{Garantir la traçabilité} de chaque prédiction (signature
        cryptographique, version du modèle, latence) pour que les
        décisions soient auditables par les régulateurs.
\end{enumerate}

% ===========================================================================
\section{Solution KOVER.IA}

\subsection{Architecture du pipeline IA}

Notre architecture orchestre trois familles de modèles complémentaires,
chacun spécialisé sur un aspect du problème~:

\begin{tabularx}{\linewidth}{@{}lXl@{}}
\toprule
\textbf{Modèle} & \textbf{Rôle} & \textbf{Fréquence} \\
\midrule
GAT (Graph Attention Network) &
Score de fraude par wallet ; capte les dépendances structurelles
dans le graphe (un wallet peut paraître propre individuellement mais
être condamné par ses voisins). &
Toutes les 5~s \\
\addlinespace
LSTM séquentiel &
Prédit la prochaine destination du flux (\emph{Uniswap},
\emph{Binance}, \emph{Hyperliquid}) à partir des cinq derniers wallets
du chemin de propagation pondéré par le score GAT. &
Toutes les 6~s \\
\addlinespace
Qwen~3 235B (Cerebras) &
Génère un rapport d'incident structuré (résumé exécutif, analyse
technique, prédiction, actions recommandées) en streaming
token-par-token. &
Toutes les 30~s \\
\bottomrule
\end{tabularx}

\medskip

Le pipeline complet ingère un flux de transactions Ethereum-like
calibré sur la distribution réelle du dataset Salam Ammari~: 83~\% de
transactions saines, 13,7~\% étiquetées \emph{Scamming}, 2,6~\%
\emph{Phishing}, plus une queue de \emph{Mixer} et \emph{Bridge}. Un
graphe roulant de 50~nœuds maximum est maintenu en mémoire sur
quatre tiers (source du hack, splitters, mixeurs, terminaux),
représentation fidèle de la topologie observée sur les hacks récents.

\subsection{Behavioral Flow Detection : la riposte aux flashloans}

Le module complémentaire \emph{kover-bfd-mev} cible une famille
d'attaques distincte mais critique~: les \textbf{flashloan attacks}.
Le sentinel se connecte au mempool Ethereum, échantillonne en continu
les transactions en attente de minage, et applique un classifieur
\emph{Direct-Drain} qui combine~:

\begin{itemize}
  \item un pré-filtre de signatures sélecteurs (méthodes connues des
        bridges Aave / dYdX / Balancer / Maker) ;
  \item une simulation \texttt{eth\_call} qui évalue le résultat de la
        transaction sans la diffuser ;
  \item un verdict IA Cerebras (~5~secondes, modèle Qwen) qui classifie
        la sévérité (\emph{LOW}, \emph{MEDIUM}, \emph{CRITICAL}) et
        l'\emph{exploit class} (\emph{direct-drain},
        \emph{oracle-manipulation}, etc.) avec un score de confiance.
\end{itemize}

Lorsqu'un \emph{halt} est déclenché, le système expose l'événement via
Server-Sent Events à un dashboard temps réel et peut, dans une
configuration production, déclencher une transaction de gel via un
multi-sig pré-déployé. Sur les benchmarks internes, le moteur traite
plus de \textbf{6,5~millions d'événements par seconde}, latence en bout
de chaîne de l'ordre de la \textbf{dizaine de millisecondes}.

\subsection{Provenance cryptographique des modèles}

Une particularité différenciante de \textsc{KOVER.IA} est l'exposition
publique d'un \textbf{manifest signé HMAC-SHA256} qui certifie la
chaîne de provenance de chaque inférence~: empreinte SHA-256 des
fichiers de poids (\texttt{gat\_model.pt}, \texttt{path\_lstm.pt}),
métadonnées d'entraînement (taille du dataset, accuracy validation,
nombre d'epochs), liste des dernières inférences avec leur latence et
leur résultat. Un auditeur externe peut télécharger le manifest, le
sérialiser de façon déterministe, recalculer le HMAC avec le secret
serveur partagé et vérifier que les prédictions n'ont pas été
falsifiées \emph{a posteriori}.

\begin{insight}
\textbf{Pourquoi c'est important.} Les régulateurs européens (notamment
sous MiCA) exigent désormais une traçabilité algorithmique pour toute
décision automatisée affectant des actifs financiers. Notre architecture
de provenance répond nativement à cette exigence, là où les solutions
\emph{black-box} concurrentes nécessitent un retro-ingénierie coûteuse.
\end{insight}

\subsection{Métriques de performance mesurées}

Les modèles ont été entraînés et évalués sur le dataset Salam Ammari
(73~034 wallets, 71~250 transactions, ratio fraude $\approx$~3,7~\%)~:

\begin{tabularx}{\linewidth}{@{}lXl@{}}
\toprule
\textbf{Métrique} & \textbf{Détail} & \textbf{Valeur} \\
\midrule
GAT --- accuracy validation & 14k wallets de test & \textbf{97,0~\%} \\
GAT --- recall sur fraude & classe minoritaire à 0,2~\% & \textbf{94,0~\%} \\
GAT --- latence inférence & batch de 50 nœuds, CPU & 9~ms \\
LSTM --- accuracy par classe & clean / Scamming / Phishing & 89~\% / 91~\% / 58~\% \\
LSTM --- latence inférence & séquence de 5 wallets & 8~ms \\
Cerebras --- débit tokens & Qwen~3 235B en streaming & 49~tokens/s \\
Pipeline complet & cycle bout-en-bout & < 5~s \\
Coverage des tests & 80 tests pytest, modules core & 95--100~\% \\
\bottomrule
\end{tabularx}

% ===========================================================================
\section{Concurrents et positionnement}

\subsection{Cartographie du marché}

Le marché de la surveillance AML on-chain est aujourd'hui dominé par
quatre acteurs, chacun positionné sur un segment légèrement distinct~:

\begin{tabularx}{\linewidth}{@{}lXl@{}}
\toprule
\textbf{Acteur} & \textbf{Positionnement principal} & \textbf{Cible} \\
\midrule
Chainalysis & Compliance \& investigations forensiques & Banques, gouvernements \\
TRM Labs & Risk scoring temps quasi-réel & Exchanges, fintechs \\
Elliptic & AML transaction monitoring & Institutions financières \\
Forta Network & Détection on-chain décentralisée & Protocoles DeFi, DAO \\
\bottomrule
\end{tabularx}

\medskip

À cela s'ajoutent des acteurs émergents spécialisés sur la détection
d'attaques (\emph{Hexagate}, \emph{Hypernative}, \emph{Ironblocks})
qui ciblent essentiellement les équipes de sécurité de protocoles.

\subsection{Comparatif fonctionnel}

Les capacités diffèrent significativement entre acteurs. Le tableau
ci-dessous synthétise les écarts les plus marqués~:

\begin{tabularx}{\linewidth}{@{}lcccc@{}}
\toprule
& \textbf{Chainalysis} & \textbf{TRM} & \textbf{Forta} & \textbf{KOVER.IA} \\
\midrule
Score IA appris (GNN)            & \texttimes & \texttimes & partiel & \checkmark \\
Prédiction destination (LSTM)    & \texttimes & \texttimes & \texttimes & \checkmark \\
Rapport en LLM streamé           & \texttimes & \texttimes & \texttimes & \checkmark \\
Manifest cryptographique signé   & \texttimes & \texttimes & \texttimes & \checkmark \\
Sentinel flashloan temps réel    & \texttimes & partiel    & \checkmark & \checkmark \\
Latence d'alerte (médiane)       & minutes    & secondes   & secondes   & < 5~s \\
Modèle économique                & enterprise & enterprise & freemium   & SaaS scale-up \\
Open source des composants       & \texttimes & \texttimes & partiel    & en partie \\
\bottomrule
\end{tabularx}

\subsection{Notre plus-value}

\textsc{KOVER.IA} se distingue sur quatre axes structurellement
défendables face aux acteurs établis~:

\subsubsection{1. La double couche AI (graphe + séquence + langage)}

Aucun concurrent ne combine \emph{aujourd'hui} un GNN (qui modélise la
structure topologique du graphe), un RNN séquentiel (qui modélise la
trajectoire) et un LLM (qui produit le rapport humain). Chacun de ces
modèles, pris isolément, existe dans la littérature académique ou
commerciale~--- mais leur orchestration en un pipeline temps réel,
avec broadcast WebSocket vers une interface unique, est notre apport.

\subsubsection{2. Provenance cryptographique native}

Le manifest signé HMAC répond à un besoin réglementaire émergent
(MiCA, AMLR-2024) que les concurrents traitent par audits externes
ponctuels. Pour un client institutionnel, la possibilité de
\emph{vérifier} chaque prédiction, plutôt que de \emph{faire confiance}
à un rapport vendor, est un argument de vente structurel.

\subsubsection{3. Latence sub-5 secondes bout-en-bout}

Là où Chainalysis Reactor produit des rapports en heures et TRM en
quelques secondes pour le scoring, notre pipeline complet (graphe~+
prédiction~+ narratif) tient sous les 5~secondes. Cette différence
n'est pas qu'esthétique~: elle déplace le cas d'usage de
l'\emph{investigation post-mortem} vers la \emph{réponse opérationnelle}.

\subsubsection{4. Surface d'intégration ouverte}

Notre interface est un WebSocket public et un endpoint REST signé.
N'importe quelle DAO, oracle, ou multi-sig peut consommer nos verdicts
sans contrat commercial préalable~--- modèle qui contraste avec les
SDK propriétaires des acteurs établis.

% ===========================================================================
\section{Modèle économique}

\subsection{Segments de clientèle}

\textsc{KOVER.IA} cible trois segments avec des besoins différenciés et
des cycles de vente distincts~:

\begin{tabularx}{\linewidth}{@{}lXl@{}}
\toprule
\textbf{Segment} & \textbf{Besoin principal} & \textbf{ARR cible / client} \\
\midrule
Protocoles DeFi (TVL > 100~M\$) & Surveillance flashloan + réponse automatique & 60--250~k\$ \\
Exchanges régulés (CEX) & Filtre AML temps réel sur les dépôts entrants & 120--500~k\$ \\
DAO de gouvernance & Alerte indépendante de la fondation, pour vote on-chain & 24--80~k\$ \\
Régulateurs / cellules TRACFIN & Capacité d'audit avec preuve cryptographique & 250~k\$+ \\
\bottomrule
\end{tabularx}

\medskip

Le segment \emph{protocoles DeFi} constitue notre point d'entrée
prioritaire~: ils ont à la fois la sensibilité au risque (un hack peut
détruire un protocole en quelques heures), un budget sécurité existant
(audit Trail of Bits, prime de bug bounty Immunefi), et un cycle de
décision court (les équipes core de 5 à 15 personnes peuvent acheter
sans passer par un comité d'achat).

\subsection{Structure tarifaire (modèle SaaS)}

Notre offre se structure en trois plans, sur le modèle classique du
\emph{land~\&~expand} avec montée en gamme par limite de débit et
profondeur d'analyse~:

\begin{tabularx}{\linewidth}{@{}lXll@{}}
\toprule
\textbf{Plan} & \textbf{Inclus} & \textbf{Tarif} & \textbf{Engagement} \\
\midrule
\textsc{Watch} &
WebSocket public, scoring GAT seul, 50 alertes/jour, dashboard partagé.
& 0 \$/mois & Aucun \\
\addlinespace
\textsc{Guard} &
Pipeline complet (GAT+LSTM+LLM), 5~000 alertes/jour, sentinel flashloan
sur 1~chaîne, support SLA~24~h, intégration Slack/Discord. &
2~500 \$/mois & 12 mois \\
\addlinespace
\textsc{Defend} &
Quotas illimités, 5~chaînes, SLA~2~h, manifest signé persisté en
BigQuery, accès au fine-tuning sur dataset client, manager dédié. &
12~000 \$/mois & 12--36 mois \\
\addlinespace
\textsc{Enterprise} &
On-premise / VPC dédié, audit OFAC sur mesure, compliance MiCA, formations
équipes, contrats négociés. &
sur~devis \linebreak (60~k\$/an+) & 24--36 mois \\
\bottomrule
\end{tabularx}

\subsection{Coûts de production}

Le \emph{COGS} d'une instance opérationnelle est dominé par trois postes~:

\begin{itemize}
  \item \textbf{Inférence Cerebras} --- tarification à l'usage, $\approx$
        7~\$ pour 1~M de tokens. Sur un pattern de production
        (un rapport / 30~s, 250~tokens / rapport, 24$\times$7), un client
        \textsc{Defend} consomme $\approx$ 720~M tokens/an, soit
        \textbf{5~040~\$/an} de coût LLM.
  \item \textbf{Hébergement backend} --- AWS m6a.large (FastAPI~+ Redis),
        $\approx$ 80~\$/mois en prod active~; 3 zones de DR pour la
        résilience~: \textbf{2~880~\$/an}.
  \item \textbf{Données on-chain} --- abonnement QuickNode \emph{Build}
        (mempool live~+ archive) à 99~\$/mois en démarrage, scalable~:
        \textbf{1~188~\$/an} par chaîne.
\end{itemize}

\medskip

Soit un \emph{COGS} consolidé d'environ \textbf{9~100~\$/an} pour servir
un client \textsc{Defend} (12~000 \$/mois~$=$ 144~k\$/an de revenu),
soit une marge brute supérieure à \textbf{93~\%}. Cette marge typique
SaaS B2B-deeptech permet de financer~: les équipes ML (recherche \&
itération sur les modèles), la R\&D produit (expansion vers Solana,
Bitcoin), et le sales/customer success (cycles de vente longs sur les
segments enterprise).

\subsection{Stratégie de financement}

Notre plan de financement projette trois phases~:

\begin{enumerate}
  \item \textbf{Pre-seed (0--12 mois) --- 350~k\$.}
        Bootstrapping sur les premiers clients design partners (3 à 5
        protocoles DeFi à 60~k\$/an chacun) plus une campagne de
        subventions Bpifrance French Tech Émergence et Horizon Europe
        \emph{EIC Accelerator}. Objectif~: closer 5 contrats annuels et
        valider la \emph{retention} > 90~\%.
  \item \textbf{Seed (12--24 mois) --- 2,5~M\$.}
        Tour de financement classique avec fonds spécialisés deeptech /
        cyber (Serena, Frst, Daphni, Inovexus) pour passer à 25 clients
        et lancer le sentinel multi-chaînes. Levée à valorisation
        préemptive de 12~M\$ sur la base d'un ARR annualisé de 1,5~M\$.
  \item \textbf{Série~A (24--36 mois) --- 10~M\$.}
        Lead investor un fonds growth crypto (Paradigm, a16z crypto,
        Variant Fund). Objectif ARR 6~M\$, expansion US (New York,
        San Francisco), équipe de 25~personnes.
\end{enumerate}

\subsection{Projection financière à 36 mois}

\begin{tabularx}{\linewidth}{@{}lcccX@{}}
\toprule
\textbf{Année} & \textbf{Clients} & \textbf{ARR} & \textbf{COGS} & \textbf{Hypothèses clés} \\
\midrule
N+1 & 8   & 480~k\$    & 50~k\$  & Design partners~+ premier contrat CEX. \\
N+2 & 28  & 2,1~M\$    & 200~k\$ & Sentinel multi-chaîne, premier régulateur. \\
N+3 & 75  & 7,8~M\$    & 800~k\$ & Lancement US, contrat OFAC, Série~A bouclée. \\
\bottomrule
\end{tabularx}

\medskip

Le break-even opérationnel est projeté à 18~mois post-seed, sur la base
d'un \emph{burn} mensuel maîtrisé (équipe technique de 6 ETP, équipe
commerciale de 2 ETP, charges fixes < 60~k\$/mois) et d'un ratio
\emph{magic number} cible de 0,9 sur les douze mois suivant la Série~A.

% ===========================================================================
\section{Conclusion}

\textsc{KOVER.IA} apporte au secteur de la surveillance AML crypto une
combinaison technique inédite~: une orchestration temps réel de trois
familles de modèles d'apprentissage profond (GAT, LSTM, LLM), une
provenance cryptographique vérifiable de bout en bout, et une latence
sub-5~secondes qui transforme l'outil d'\emph{aide à l'investigation}
en \emph{capteur opérationnel} directement intégrable dans les workflows
de réponse à incident des protocoles, des exchanges et des DAO.

Le projet, livré dans le cadre du Hackathon AI Paris~2026, est déjà
fonctionnel sur l'ensemble du pipeline~: les modèles sont entraînés sur
des données Ethereum réelles, le manifest cryptographique est exposé,
le sentinel flashloan détecte les attaques en temps réel, et l'intégration
continue (GitHub Actions~+ SonarCloud~+ 80 tests pytest avec une
couverture supérieure à 95~\% sur les modules métier critiques) garantit
la fiabilité de la base de code. La prochaine étape consiste à valider la
proposition de valeur auprès des design partners identifiés, à étendre
le sentinel à Solana et BNB Chain, et à boucler la levée seed pour
constituer une équipe de produit complète.

\begin{quote_kv}
La fenêtre d'opportunité est étroite~: la régulation MiCA entre
pleinement en vigueur fin 2026, et les protocoles DeFi qui n'auront pas
adopté un système d'AML traçable et auditable seront mécaniquement
exclus des marchés européens. \textsc{KOVER.IA} se positionne comme
l'option \emph{native crypto} pour répondre à cette contrainte, là où
les acteurs traditionnels (Chainalysis, Elliptic) restent ancrés dans
un modèle de \emph{compliance lourde} hérité de la finance bancaire.
\end{quote_kv}

\vfill

\begin{insight}
\textbf{Annexes disponibles.} Documentation technique complète sur le
dépôt GitHub~: architecture du pipeline, scripts d'entraînement,
manifest signé, jeu de tests pytest, configuration SonarCloud. Les
modèles entraînés (\texttt{gat\_model.pt} et \texttt{path\_lstm.pt})
sont versionnés dans le dépôt avec leur empreinte SHA-256 vérifiable
sur l'endpoint \texttt{/ai/proof}.
\end{insight}

\end{document}
